%%%%%%%%%%%%%%%%
% SPACE EMPIRICS
%%%%%%%%%%%%%%%%
\section{Space empirics \label{sec:spaceempirics}}

The job of this section is to provide an overview of the reality of space costs. This is not meant as a substitute for empirical data, but instead to ward off the dangers of going too far in the direction of \emph{a priori} model-building, where `common sense' assumptions about space can easily be shown to be contradicted by the briefest of glances at empirical reality. Chapter \ref{ABM or another one?} explores in more depth how the modeller can approach this trade-off between model and real world: to pre-empt that discussion, all models must contradict reality to some extent. It is the modeller's job to be aware of where that happens, and so be able to correctly interpret the model given its particular limitations.

Regardless of how simple a model is, as full an understanding of the complexities of reality as possible will make model interpretation better. A particular target for critics of location theory has been the use of abstract Euclidean space: a featureless plane, clearly nothing like reality. One job of this section is to examine this particular assumption, explain why it is being kept. The conclusion will then come back to this as an example of thinking about the trade-off between model simplification and a more complex reality. (Was the simplification justified? How to tell if it was justified? Noting that 'because is was easier / simpler' cannot justify it by itself...)

\subsection{Terminology \label{sec:spaceterms}}

\begin{itemize}
	\item Space costs: a top-level term used to describe any costs connected to space and location, including all further terms below
	\item Distance costs: incurred as a result of the physical distance between places
	\item Land costs: ... already written about somewhere, need to bring in here
	\item Density costs: relates to congestion charges; used here as proxy for land costs, see section \ref{sec:densitycostsetup}.
	\item Trade costs: a catch-all term for all costs incurred between trading partners; see section x.
\end{itemize}

That is potentially quite a confusing list, but it is important to state it explicitly, and work through the theoretical inter-relations between the terms. 

The definition of a `cost' itself, economically, needs some explanation. One of the central goals of economics is to understand the impact of \emph{all} costs. In practice, what this means is that often a mathematical description of costs will be used to discuss quite a range of real-world cost equivalents. 

%Need to x-ref forward with this, I think. And, is this the place for this discussion? It could perhaps go in the philosophy section? Which is then starting to look like a separate chapter again.

%Think I want to say here: you end up with e.g. something like density cost standing in as a proxy for several other costs - and that's OK. 
%Also cx2 2nd nature cost... 

The quantity of fuel needed to move a given distance is one of the few objective elements involved.  The total cost is a much messier proposition. 

%%%%%%%%%%%%%%%%%%%%%%%%%
% Measuring space
%%%%%%%%%%%%%%%%%%%%%%%%%
\subsection{measuring versus modelling space \label{sec:measuringspace}}

Chapman: Good line: `because of the nature of freight rates, cost-space and absolute space are rarely isomorphic.' [179/15]

%Ontology/epistemology

It would be nice if there was a clear way to separate the reality of space costs from how they are theorised. Things are messier than this, however: they cannot be boiled down simply to what can be measured. An example where measurement is a particularly thorny problem can be found in the trade cost literature; as Anderson and van Wincoop (\citeyear{anderson_trade_2004} p.692) point out, ``a theoretical approach is inevitable to infer the large proportion of trade costs that cannot be directly measured in the data''.  

Such an approach has already meshed what objective component of space cost exists with a theory about how actors will respond to them. 

Most common method for describing space cost in aggregate: distance decay functions.  Gravity models.  Chapman 180/15.  High decay goods

Anderson and Wincoop come up with numbers describing the `tax equivalent' costs of trade for industrialised countries: 170\%. This breaks down into ``21\% transport costs, 44\% border related trade barriers and 55\% local retail and distribution costs.'' [626]
%a place where the reliance on `empirical' answers is very model-heavy. 

They use a gravity model, though one they consider superior economically, because theoretically grounded. Section \ref{sec:micromacro} discusses the example of gravity models and theory in more depth; here, the goal is trying to identify the reality of space costs. So what actually lies underneath the aggregate percentages Anderson and Wincoop have produced? It is difficult to say, since there is such a large gap between that reality and the figures, via their particular methodology.
%This would be an argument for building models to help with ontological prediction: to know what to look for. 

It is not an easy problem to fix. Caught between two different issues: it is possible to get data for basic transport costs - that is, the actual money cost of moving stock across space, as a proportion of overall business costs. But as the trade cost puzzle illustrates, this approach clearly misses a large chunk of space costs. 

Firms primarily want to locate near to labour: this is a key finding from the literature. It can be understood in space-cost terms by comparing with the fact that service providers are often reliant on proximity to customers. In both cases, the same \emph{type} of space cost exists: that incurred by people moving location. 


Evidence shows that in many countries, distance of trade has actually been decreasing - ``distance impedes trade by 37\% more since 1990 than it did from 1870 to 1969'' (\citealt{carrere_distance_2009} p.6) This affects poorer, more distant countries much more than the main centres of economic activity. 

There is an empirical question here, overlapping with studies examining the trade cost problem: how to find out, firstly, what has actually been happening and secondly, why? Role of theory.

DOT drops as EU / US subcontract to Eastern Europe and S. America. [33/20] political developments, then, not changes in space costs.
cf. other political points: Whole issue is a difficult problem for abstract analysis.
European Union and U.S. Retail Gas taxes: 60 to 75 per cent and 20 to 25\% respectively.
Road tax; Motor Carrier Act etc. [23]
Should be put next to Fujita's points... where are they now?

Some types of space cost data are relatively easy to acquire, such as tariff data. Others, as Behrens points out, are so challenging that workarounds are needed (\citealt{behrens_changes_2007} p.626). 

Behrens notes that \gls{ge}  presumes ``the different components of trade costs can be reduced without loss of generality to a single parameter'' [626] and that ``the underlying assumption is that all the components of trade costs affect industry location and trade patterns in the same way.'' \citealt{behrens_changes_2007 p.626}

%%%%%%%%%%%%%%%%%%%
% Route development
%%%%%%%%%%%%%%%%%%%
\subsection{Route development}

Route and infrastructure development are by far the biggest modelling challenge: it is just as important as basic fuel costs and technological change, but it is the most byzantine. The largest long-run changes are due almost entirely to infrastructure development; most obviously, choices about new rail and road routes. Andersson and Stromquist make a particularly strong claim: ``all the major transitions in the European economic systems were accompanied, or initiated by major changes in transport and communications infrastructure'' (\citealt{batten_transportation_1988} quoted in \citealt{button_transport_2010} p.421).

So infrastructure \emph{change} is pivotal: it changes proximity unevenly. Transport development tends to happen between economically important sites and, as Fowler notes, ``uneven power relations'' are both a cause and consequence (\citealt{fowler_reexploring_2006} p.1433). Routes themselves are also part of feedbacks: it is, as Chapman says, ``equally plausible to regard transport investment as a result of a need for movement or as a generator of movement'' (\citealt{chapman_people_1979} p.230).

Route and infrastructure development, then, is a major modelling headache for space costs. For a start, it makes no sense in anything other than $R^2$. Between or within discrete regions, it must be conceptualised as either an externality or a generic transport cost change between points. In $R^1$, the same applies: there is only one route between points. Choosing discrete space or $R^1$, however, does not make the problem go away: it just masks the fact of this uneven change dynamic. If, as argued here, it is perhaps the most important long-run dynamic, this means avoiding it leads to wrong results. In $R^2$, however - as on the locally two-dimensional reality of the Earth's surface - this change can at least in theory be examined.

This thesis wants to argue that 'Weber-type stories' absent from \gls{ge} are worth telling, but route and infrastructure development makes telling them very difficult. $R^2$ has a uniquely complex relationship to long-run space cost change. How could one possibly go about modelling it in a way that fits with \gls{ge}'s goal of endogenising outcomes? The obvious answer is to avoid it: as long as routes in $R^2$ are presumed to be static, it is possible to assume a mapping of any route system to a Euclidean approximation, making the analysis potentially much easier. Specifically, it becomes possible to use Euclidean distance as a reasonably accurate estimation of actual distance and space cost between points (see, for instance, \citealt{cooper_use_1983}.) This is only in any way plausible, though, for considering short to medium-term change: time-scales where the impact of route development will be only minimally felt. It is also quite a large assumption that fixed networks can be mapped onto a Euclidean plane unproblematically: see, for example, Wilhite's work on fixed networks (\citealt{wilhite_economic_2006}). Nevetheless, this approach has one clear advantage: it is very simple. Calculations are trivial and, as an added bonus, visualisation is straightforward and intuitively easy to grasp. This method is used in this thesis; the rest of this subsection, however, explores the issue, to give this approach some context.
%Out of all model assumptions, it is this one - not utility or other economic simplifications - that makes it most questionable.

%There's more of a story about short, medium, long term spatial change and e.g. knock-on effects of energy costs, energy type, mode etc. May relate back to transition stuff. 

A simple way to describe the discrepancy between route distance and straight line distance is with a `route factor' (see e.g. \citealt{chapman_people_1979} p.215; \citealt{black_transportation:_2003} p.68). This is simply the ratio of two distances: the distance between points following a network route, over the Euclidean distance between those points. This value approaches one as the two measurements converge. Chapman notes an early study by Timbers (\citeyear{timbers_route_1967}) that examined routes between 780 pairs of towns: the average route factor was 1.17. As places develop, then, networks move towards complete connectivity: there is a well-established correlation between network density and development levels (\citealt{chapman_people_1979} p.220). This suggests that the `assume short-to-medium term space cost change' approach will be less applicable where infrastucture is less developed.
%oooo, in the wrong order still!

In the seventies, Haggett \emp{et al}, after an exhaustive look at the issue, concluded that ``there remains to be specified a comprehensive model of route development'' (\citealt{haggett_locational_1977} p.95). No such model exists today, though perhaps for good reason: Black muses, ``perhaps it is best to have a field whose members have skills that can be applied to contemporary problems, even if such applications slow the field's conceptual and theoretical development'' (\citealt{black_transportation:_2003} p.9). Network analysis itself is, of course, now a very large part of the economics literature (for an overview, see e.g. \citealt{goyal_connections:_2009}) and has a central role in the development of complexity approaches to economics (see section \ref{} in chapter \ref{} for more detail on this point). There is also network-related discussion in the \gls{ge} literature, but again, its focus is on \emph{economic}, not transport, networks; for example, see \citealt{johansson_agglomeration_2003}). 

More abstract network analysis relevant to transport optimisation has found an enthusiastic research community in physicists and statisticians. A particularly good recent example of this is Gastner and Newman's network growth model (\citealt{gastner_shape_2006}), in which they explicitly look for the balance between shortest routes and total route distance (the more total route, the more expensive the build), concluding that while ``these two criteria are often at odds with each other... real networks nonetheless manage to find solutions to the distribution problem that come remarkably close to being optimal in both senses'' and that `growing' them in their model does as well, or better, as attempts to explicitly optimise them. 
%The two ideal solutions were ``one which minimises the user's travel costs and one which minimises the builder's construction costs.'' [218/16] What is 'compromise' from user/builder tension?

This type of analysis, however, tends to be economics-free: cost is reduced to route length alone, and route development is not driven by particular interests. For instance, Gastner and Newman's distinction between shortest route and least total route optimisation matches the difference between optimal networks for users versus builders (\citealt{haggett_locational_1977} p.218). Users want dense networks to lower costs of travelling between points; builders want as few links as possible to reduce their costs. Once development, routes are also ``subject to international arbitrage by profit-maximising agents'' (\citealt{behrens_changes_2007} p.626)\footnote{This is another example of a point that is made in human geography textbooks from the seventies or before, e.g. \citealt{chapman_people_1979} p.117: ``the principle of charging `what the commodity will bear' is widely adopted by carriers and this ensures that high-value goods, for which the cost of transport will add proportionately less to total costs, usually pay more to move a given distance whereas goods of low unit value may be charged rates which do not even cover the costs involved.''}. Given the nature of transport routes, this can often mean monopoly, monopolistic competition and cartels, as with (until recently) shipping 'Conferences' (\citealt{carrere_distance_2009} p.17), American trucking prior to the Motor Carrier Act (\citealt{glaeser_cities_2004} p.204), or - an example from a particular industrial sector - `basing points' for steel prices in Europe and `Pittsburgh plus' rules in the U.S (\citealt{chapman_people_1979} p.117).

cf. Ladley... but then Wilhite above. Eek. 

Pose the question instead as: how do incentives of different actors interact with cost changes to produce route and infrastucture morphology? It is, first, a very important and largely unsolved problem, and, second, beyond the scope of this thesis. But it must not be forgotten when discussing the results of any model wanting to say something about space cost change in two-dimensional space.
%Is a drive for endogeneity useful? Is it the same thing as looking for generality?
%This sort of analysis is much less a part of geography now, tending instead to be grounded in the more recent physics-led resurgence in network analysis. A particularly good example of this is 

Returning to Haggett's pointing out the lack of a comprehensive route model: one might argue it is simply impossible to develop such a general, endogenous model that could parallel the drive for endogeneity in other areas of economics. Gastner and Newman's work, for example, assumes given network vertices and allows the growth of edges from those, and network economics tends to consider action on fixed networks. Behrens (\citeyear{behrens_changes_2007}) take a model `case study' approach to see how a small fixed network with intervening opportunities effects a space cost change situation. 
%More interesting: what is the outcome from these tensions? Who actually drives route choice? It isn't a simple reflection of consumer demand. 
%precisely the problem: normative analyses in transport geography \emph{are} connected to reality - perhaps more than any other economic realm, transport is a matter of public policy, not merely the result of competing private interests - though of course, reality is some fiendish tangle of these.


Impact of existing `colonial' transport networks e.g. S America. Argentina, Uruguay: railways for shipping meat, wool to Britain. Chile, Bolivia, Peru, for copper, lead, tin to Pacific coast. Gilbert 1974: Latin American Development. London. [231]
Commercial relations between countries thus limited. (Hmm - what role would this play in overall DOT stuff, I wonder? And how people try and determine why trade flows are as they are�)
``Latin America appears as a complex of national space more closely tied to exogenous decision-making centres than to itself.'' Quote from Melchior 1972, `The integration of space in latin america', Latin American Urban Research, 2, p.88. [16]

So compare to DOT: what role route development over time, and the interests? Hugely interesting: still true now that no model exists; is that because no generalisations are possible? At any rate, biggest area that I can't do here.
%Note comparison to MAR and Jacobs externalities in that paper though... 
%May want to go back and check the WoS citations for that again, if possible: some interesting stuff. And plenty of time, eh?

%%%%%%%%%%%%%%%%%%%%%%
% OTHER TRANSPORT BITS
%%%%%%%%%%%%%%%%%%%%%%
\subsection{Other transport bits}

Concentrating on network formation masks a further complexity: what actually happens on the network. Here are two ways this can be thought about. The most obvious is to consider the cost of using space, as reflected in the basic facts of congestion. Glaeser and Kohlhase describe the problem (for the U.S. again): 

\begin{quote}
``Over the past three decades, congestion and delay have been increasing in all size classes of cities, not just the very large urban areas.  For all and 75 urbanised areas taken together, on average the annual delay increased by over 280\%, with the most dramatic increase occurring for large cities - those of between one and three million -about 450 per cent.'' (\citealt{glaeser_cities_2004} p.209)
\end{quote}

%An example used in this thesis, from the list above, is density cost: as section \ref{sec:densitycostsetup} explains, a cost based on density of actors can stand in as proxy for combined land and congestion costs. 
%As with congestion, many of these costs can often be difficult to quantify, though 
%"one between 1980 than 2000, commute times rose by about 13% for the combined 75 areas." 
% (though not impossible, see \fixref{congestioncostrefs}), but this fact does not rule out their inclusion in a modelling problem. 

Another way to think about how space is being used economically is `transport intensity', defined as ``the ratio of gross mass movement to gross domestic product (GDP)'' peak 1994 in Scottish report 4/12
Interesting: if one were to suggest that transport costs were dropping because good value was increasing, then transport intensity should not be increasing. It is. 

Falkner: ``Official statistics show that from the time of the first oil shock in 1973, to today, domestic households, industry and services have been able to reduce their reliance on oil products. Not so in the case of transport.'' [28/17]

question: are there any spatial structures with LoV where localised transport intensity goes up...?

\subsection{how people actually cost space / does it matter?}



While some argue 'it is better to assume that moving goods is essentially costless' (Glaeser) because they are so low, others point out that 'our basic local needs have become dependent upon a hyper-complex, integrated, tightly-coupled global fabric of exchange' (Korowicz) 

humpty dumpty!

`Glaeser hypothesis' transport costs are so low that they need not be considered in analyses. Other space costs are still important for Glaeser, though - particularly those that affect how people can use their time. 


Glaeser and Kohlhase note both the money and non-money movement costs people face: they cite a 2001 consumer expenditure survey showing that, in the U.S., ``18\% of total expenditures for the average household is spend on vehicle purchases, gasoline and other vehicular expenses'' - but point out that time costs are not accounted for in that, and that ``these time costs are not withering away with technological progress'' (\citealt{glaeser_cities_2004} p.208).



There is a hard limit to people's time: a flow of twenty-four hours per day must be spent. The spatial consequences of this are most strongly felt for those activities requiring the most time - and employment tops this list. Close seconds are access to housing and goods. The spatial solution to this problem does not necessarily require people to stay within a day's limit: Ojo, for example, studied agricultural `commuters' in 1970s Yorubaland who developed a ``fortnightly pattern of periodic commuting to reduce the friction of distance'', moving between more distant settlements (\citealt{ojo_journey_1973} p.86; \citealt{chapman_people_1979} p.122). Closer to home, commute patterns that span larger distances than one day's range do exist - but the vast majority keep themselves to the geographical limit bounded by that one day. 
%Fascinating conclusion to that paper: I can see the whole seeing like a state thing going on there. I wonder what sort of schemes did in fact get carried out? Begin to see why people would object to atempts by geographers to make things more efficient... I think this paper should be used as an example of some of the problems with modelling and what modellers should be doing with them: there's a world out there, beyond the model itself or the data collection!
%http://nigerianwiki.com/wiki/Afolabi_Ojo - fascinating!

The idea of external costs comes in particularly useful here. There is a cost to firms derived from people's movement, whether those people are travelling to work or going to get a haircut. In the latter case,  the service provider has externalised one of the major costs: getting the person�s head under the scissors. In the case of commuting, the cost is less cleanly external or internal to the firm's production structure, since wages in part must cover it (and one might argue that they completely cover it, if wages are determined by marginal productivity.)
%really?

Might rent, in effect, internalise these space costs though?
%There's a thunen argument here somewhere... 
(Model scenario \ref{sec:firmsLoV} has an example of such an external cost and how two firms cope with them.)

%Button: "after housing and food, transport represents, at about 10 to 15 percent for a single-car-owning household, the largest element of expenditure out of income in industrialised countries." [90] tie to theory of people as productive units.
%Button: "Time savings at work permit a greater output to be produced within a given time period which, again, drawing on the marginal productivity theory of wage determination, will be reflected in the marginal wages paid." [10/100]

Leamer/storper: Cost of human containers rising - congestion, real wages

``humans remain the containers for shipping complex uncodifiable information. The time costs of shipping these containers is on the rise because of congestion on the roads and in the airports while the financial costs of so doing are also rising due to increases in real wages of knowledge workers who are the human containers'' (\citealt{leamer_economy_2001} p.648).

"Over the past three decades, congestion and delay have been increasing in all size classes of cities, not just the very large urban areas.  For all 75 urbanised areas taken together, on average the annual delay increased by over 280\%, with the most dramatic increase occurring for large cities - those of between one and three million -about 450 per cent." [209]
"one between 1980 than 2000, commute times rose by about 13\% for the combined 75 areas."  [23]

Chapman 112/13: tapering rates for distance cost.  Reflects fixed vs variable costs. Cf. dwell costs [19]
Blanket rates, zonal: because more specific charging difficult to set up, communicate to customers [14]
Mechanism design to externalise cost: milk marketing board 14
Chapman again: principle of charging what the commodity will bear.  Cost of space subject to supply and demand.
Break of bulk charges at ports; Pittsburgh plus compare to basing point for steel production
Cartels etc. shipping 'Conferences' [dot17/20]
Ah, this is a point made for JIT: more important now. "the quality and reliability of a transport service are often more important to the industrialist than the payments made for it." Ref, bayliss & edwards '68; wallace 74. [chapman 117/15] time factor can be critical
�
Tech changes act as gateway: e.g. Chapman: refrigeration for meat opens up meat export for Argentina and New Zealand 182/15

and make decisions: 

button: acceptable sites chosen over a protracted and thorough search [8]
G/K: "using the same workers is approximately eight times more likely to increase the degree of co-location than are trade relationships." [223] 

Chapman: smaller firms less economically led; large firms not short term profit but long-term stability and growth.


\subsection{proportions}


Adapted landscapes: Glaeser hypothesis: transport cost can be ignored.  But landscape reflects adaptation to those costs.
The fundamentals do show themselves: transport costs may be negligible, but that's only because the economic landscape has already adapted itself to minimise them, into a series of overlapping plates of economic activity. This goes back to the problem of analysing the complete impact of costs, not just those that can be measured - a point made by Burstein \emph{et al}, who note that concentrating just on incurred costs is - 

`` - likely to underestimate true transportation costs. This is because we do not observe the cases in which transportation costs were so high that the transaction did not take place.'' (\citealt{burstein_distribution_2003} p.1198)

This is not a new idea: `demand thresholds' mark where high costs cut off transactions (\citealt{chapman_people_1979} p.86), but it has tended to be applied to consumer demand problems, rather than the economy as a whole. It is certainly easiest to understand Burstein's point by applying it to people: as discussed above, time thresholds mark a clear limit on how far people can travel to work. The same applies to Firm inputs other than labour, however. While the problem is easy to state, it is harder to work out how those costs could be quantified. After all, it is essentially a counterfactual: what trades would have taken place if costs were lower? Luckily, it is easy to ask a model to do this comparison.

This point overlaps with the `ecological criticism' discussed in section \ref{sec:ecocrit}: it appears, mathematically, that transport costs can be reduced simply by \emp{increasing} the cost of other inputs into the production process. It appears that simply by making other inputs more expensive, the importance of space can be lessened, with no limit on how much they can be reduced. 

A spatial version the ecological criticism might go like this: it is obviously nonsense to think that space can be made unimportant simply by increasing the cost of other inputs. There are hard cost limits to moving goods and people. It is a complicated argument to unpick: it begs questions about, firstly, where value comes from in the first place and, secondly, where those extra non-space costs might originate.

However, it does appear to be the case that space costs \emph{relative to} other costs are very important. Consider the service industry: Glaeser and Kohlhase correctly point out that since ``transportation costs for people are still high, but transportation costs for goods are not'' (\citealt{glaeser_cities_2004} p.220), proximity to people is the most important spatial factor for services. They also note (p.201) that ``services tend to involve little freight shipment''. But what about transport costs as a proportion of the service sectors' overall costs? For the U.K. at least, Diamond and Spence found they were 9.9\%, compared to 4.7\% in manufacturing (\citealt{diamond_infrastructure_1989} in \citealt{button_transport_2010} p.58). Unsurprisingly, much of the service sector has lower overheads than other sectors. The paradoxical result is that services are more space cost sensitive. 

Of course, the space cost of all firms' biggest input takes priority: proximity to the labour market. The labour is also the largest source of extra value input into the production process. The majority of value embodied in transported goods comes from this source, rather than the intrinsic value of the good itself. 
\fixref{any citations for labour as largest source of value?}

As section x outlines in full, a simple way to think about this is to consider the unit price for a unit of a good at a given distance. If $P$ is the full unit cost, $p$ the base per unit cost, $c$ the cost to move one unit of good over one unit of space, and $d$ the distance moved, then:

\begin{equation}
P = p + (c\cdot d)
\label{eq:spacecost}
\end{equation}

$c$ can stand in as an index for the generic costs of moving goods, whether due to weight, bulk or some other factor. (It could also be defined to include other factors that produced increasing returns to distance; that is, the further a unit is moved, the cheaper per unit it becomes, as variable delivery costs increase relative to fixed delivery costs.) For the purpose of this section, the important point is to note the impact of increasing $p$ - for example, if a good had more skilled labour input, thus increasing its base value. The full unit cost $P$ would become less dependent on distance.

One way of thinking about this effect is to use the idea of `value density'. This can be defined as the ratio of product value to physical size or weight (\citealt{lovell_product_2005} p.144); here, that is $\frac{p}{c}$. There are two particularly important empirical considerations for value density. 


in single firm, opt, which illustrates the spatial impact of value density: the more value is added along supply networks, the more feasible large space costs become.

A large input of skilled labour tends to happen in firms with bigger economies of scale. Lovell \emph{et al} note Coopers' research on the production of microchips as a classic example of combining massive scale with very high value density (\citealt{cooper_logistics_1990} in \citealt{lovell_product_2005} p.144). Given the definition of increasing returns in production from section \ref{sec:prodfuncts}, this is unsurprising: it is exactly that a unit of input labour can produce a higher output. That output, however, is measured more in value, rather than any increase in physical good, the more skilled the labour input involved. As Lovell \emph{et al} note, there are obvious space cost consequences: ``as the value density of a product decreases the percentage of additional distribution costs increases.'' (ibid.) 
%They have an example too that's probably worth going back for a look at.

%Examples of value density-ish stuff: wood products av shipment length of 287 miles - if by truck, transport costs  =1/5 value. Base metal just under 1/10th (if by truck.) Commodity dictates mode? Thus geography, at least in part?
%Cf. machinary, electrical: always less than 1.2\% transport costs. [207] According to their analysis: only 18\% of all shipments occur in industries where transport costs are more than 6\% of total value - even if all transport was by truck."  Q: how is JIT affected by mode choice? [23]

%Value density dictates mode.

%This leads to a somewhat odd conclusion: space costs could be reduced simply by making people 
%how people value something can determine spatial morphology. (cf. Becker's claim that changing values make no difference - they probably do spatially.)
%Can increase value density as well as lower transport costs.

Chapman: 182/15. UK primary, secondary production: transport cost = 3\% in 1968. Dawson 1977: `costs of production and location of british industry.' Geography 62, 93-6. Highest for bulky, low-value materials: still only 20\% of costs.
G/K: ``in a service economy where transport costs are small and natural productive resources nearly irrelevant, weather and government stand as the features which should increasingly determine the location of people.'' [212]
Compare to Button: service sector more transport intensity
�
G/K: Transport costs: 9\% to 2\% of national product (in the US presumably) 1900ish, now. [22]
``a significant fraction of GDP is not shipped. Services tend to involve little freight shipment.'' [201/23] But they are very location and transport cost sensitive, relatively. Discuss.

``Overall, across all modes there have been declining costs, and in the absence of outside factors (the oil crisis, government regulation) costs per ton-mile, within each mode, appear to be declining by about 2\% a year.'' [204/23]
Relative to overall GDP.  Which of course means a constant upward trend in value density.  Note ``secular trend in value of dollars per shipment.'' [23]

Could mention Isard here? Coal and steel...


\subsection{risk is expensive}

Risk is a major cost for firms, though a complex one to theorise. It is important enough for real-world outcomes that it is worth discussing briefly. Note, however, that while uncertainty does play a part, explicit risk modelling is not part of this thesis (though it is a highly interesting topic for agent modelling).

Food: Chapman.  Farmers cultivating many distant plots hedge bets.

So. Chapman [23] General findings: smaller firms less economically led; large firms, not short term profit but long-term stability and growth. Hmm.
So, larger firms always about managing risk to stay stable over time; more-so now, lean retailing etc...

While more recent research highlights the risk-reducing properties of the modern trans-national corporation, the importance of risk-aversion in transport is not new. Wallace, for example, found in the U.K. that ``where the choice between two modes involves both a considerable price and a considerable quality-of-service differential there is a clear willingness to pay the necessary price of reliability'' (\citealt{wallace_relationship_1974} p.41). Certain goods have always been more time-sensitive than others. Chapman also notes that ``consignments which fail to turn up on time impose indirect costs by interrupting production schedules. Indeed, in certain industries, the time factor is critical'' (\citealt{chapman_people_1979} p.117).

Oil risk volatility is very damaging: "not a simple dependence of the economy on the level of oil prices, as would be suggested by production-function based accounts." [Hamilton 33/22] Disruption! Dislocation! Etc!
"small mismatches in energy supply and demand produce wide price swings." [16] 
This interpretation is supported by the fact that investment decisions are delayed both when oil prices increase and decrease. 


Just-in-time production.  (Intra firm only?)

Oil Crunch: Volality is very damaging. This is also true for almost any economic feature: risk is expensive.


One of the solutions to deal with uncertainty is the multinational firm: it combines scale and, as Leamer and Storper point out:

\begin{quote}
``much global trade consists merely of shipping products or components between divisions of the same firm located in different countries - transactions that do not raise the trust and enforceability issues present in arm's length transactions.'' (\citealt{leamer_economy_2001} p.649)
\end{quote}

Because of this, modern optimisation techniques for global production applied in-firm have been very successful in managing spatial risk. Carrere \emph{et al} (discussing \citealt{abernathy_stitch_1999}) note that for apparel ``the key to success is no longer solely price competition but the ability to introduce sophisticated information links, forecasting capabilities and management systems'' (\citealt{carrere_distance_2009} p.34). This makes increasingly `lean retailing' possible, and has led to other geographical effects. Particularly, these lean firms have been bringing production networks closer: U.S. firms networked with South America and European firms with Eastern Europe.

G/K: "manufacturing has declined as well and now represents only 14.1\% of the economy."
This is all U.S.  What about outsourcing of production globally?  What impact on space costs?  Has there been any actual drop in weight/distance of consumed goods?

\subsection{change?}

Efficiency	Different impact of higher costs on supply and demand. Cf.  Research and development on producing more efficient cars to people's car purchase decisions shifting to already extant efficient vehicles.

Relation of change to model choice. This probably shouldn't be here, but -

Transport economics: short-term, button: taking a "given land-use pattern and looks at methods of providing efficient transport services within this constraint" [8/51]
Button again: supply chain analysis in management literature; operations research etc [11/327] 





%%%%%%%%%%%%%
% CUTTINZZZZZ
%%%%%%%%%%%%%

%Last-mile problem.  Mentioned in button [11/341]
%Marchand 1973: Venezuela using oil to improve road network. 1936-61 = av travel time between nodes 3000 to 700 hours. [chapman7/13]
%For comparator intervening opportunities: local development of routes may actually meet some of the fixed costs of the total route contributing to development of routes.  Chapman: 183/15
%The reason it is impossible to maintain a clear distinction between ontology and epistemology of space costs is of course the complicating human factor.  A theory is required...
%Analysis of infrastructure change does exist; operationally, it needs to. Current plans for high-speed rail routes in the UK, for example, have been analysed using an econometric framework developed at the London School of Economics (\fixref{overman report}) for assessing the economic impact of transport network choices. Of course, this has to make an assessment of the existing economic landscape. However, to find more theoretical attempts to understand the cumulative effects of long-term change (`history of the world'-type questions mentioned above) one has to look to pre-1980 geographers like Haggett \emph{et al} and Chapman (\citealt{haggett_locational_1977}, \citealt{chapman_people_1979}).
%This presents something of a catch twenty-two. 
%Black also highlights the importance of Haggett \emp{et al} in embedding network analysis etc... p.7
%Cf. possible outsourcing of production intensity. Certainly not the case that we're consuming less, is it?
%Deardoff 1984 p.470: ``like frictions in physics, transport costs are almost universally ignored in trade models in the sanguine hope that if included they would not materially affect the results.'' In behrens 626/18
%``The transport sector accounted for some 80\% of the demand growth between 1990 and 2000, and became the largest demand-side sector.'' Button [10/187]
Scotland report suggests this is way of measuring decoupling.  But cf. Glaeser: is it decoupling actually driven by outsourcing?
Is reduction in energy and emissions easier and volume? 4-5/13
Three components of intensity: traffic(distance), energy carbon
Leeds project: little evidence to suggest that the growth of transport demand has been or ever can be decoupled from the rate of economic growth. 7/13
love-of-variety based trade models imply intensive horizontal intra-industry trade

%One can consider households as productive units; main input into production: labour.  But sort of externalised.  If treating households as productive units getting food counts as production cost. [9]
%Chapman: journey to work in agriculture or urban reversed.

